\chapter{การเรียนรู้แบบไม่มีผู้สอน}
\label{CH6_Unsupervised}

การเรียนรู้แบบไม่มีผู้สอน (Unsupervised Learning) เป็นสาขาหนึ่งของการเรียนรู้ของเครื่อง (Machine Learning) ที่มีความท้าทายและน่าสนใจไม่น้อยไปกว่าการเรียนรู้แบบมีผู้สอน (Supervised Learning) ที่ได้กล่าวผ่านมาแล้วในบทก่อนหน้า ในการเรียนรู้แบบมีผู้สอน ตัวแบบจะได้รับข้อมูลที่มีการติดป้ายกำกับ (Labeled Data) พร้อมกับคำตอบที่ถูกต้อง (Ground Truth) ซึ่งทำให้สามารถวัดความผิดพลาดและปรับแต่งตัวแบบได้โดยตรง ต่างจากการเรียนรู้แบบไม่มีผู้สอนที่ตัวแบบต้องค้นหาความรู้และรูปแบบ (Patterns) จากข้อมูลที่ไม่มีการติดป้ายกำกับใด ๆ ทั้งสิ้น ลักษณะเด่นนี้ทำให้การเรียนรู้แบบไม่มีผู้สอนเป็นเครื่องมือสำคัญยิ่งในยุคที่ข้อมูลมีปริมาณมหาศาลแต่มีการติดป้ายกำกับ (Labeling) ที่มีราคาแพงหรือใช้เวลามาก

เนื้อหาในบทนี้จะนำเสนอหลักการพื้นฐานของการเรียนรู้แบบไม่มีผู้สอนอย่างเป็นระบบ โดยแบ่งเป็น 4 หัวข้อหลัก ได้แก่ \textbf{หลักการและความแตกต่างจากการเรียนรู้แบบมีผู้สอน} เพื่อทำความเข้าใจพื้นฐานทางทฤษฎี จากนั้นจะศึกษา \textbf{การจัดกลุ่มข้อมูล (Clustering)} ซึ่งเป็นเทคนิคที่ใช้กันอย่างแพร่หลายในการแบ่งกลุ่มข้อมูลที่มีความคล้ายคลึงกัน และ \textbf{การลดมิติข้อมูล (Dimensionality Reduction)} ที่ช่วยจัดการกับข้อมูลที่มีมิติสูง รวมถึง \textbf{การตรวจจับความผิดปกติ (Anomaly Detection)} ที่มีความสำคัญอย่างยิ่งในงานด้านการรักษาความปลอดภัยและการเฝ้าระวังระบบ เทคนิคเหล่านี้ล้วนเป็นเครื่องมือที่จำเป็นสำหรับนักวิทยาศาสตร์ข้อมูลและวิศวกรที่ต้องทำงานกับข้อมูลในโลกจริง

%####################################
\section{หลักการและความแตกต่างจากการเรียนรู้แบบมีผู้สอน}

การเรียนรู้แบบไม่มีผู้สอนมีปรัชญาการเรียนรู้ที่แตกต่างจากการเรียนรู้แบบมีผู้สอนอย่างสิ้นเชิง ในการเรียนรู้แบบมีผู้สอน ตัวแบบจะได้รับข้อมูลนำเข้า $\mathbf{x}$ ควบคู่กับค่าเป้าหมาย $y$ ซึ่งเป็นคำตอบที่ถูกต้อง และเป้าหมายคือการเรียนรู้การจับคู่ระหว่าง $\mathbf{x}$ กับ $y$ ให้ได้อย่างแม่นยำ ตัวอย่างเช่น ในการจำแนกประเภทผู้ป่วยโรคเบาหวาน ตัวแบบจะได้รับข้อมูลระดับน้ำตาลในเลือด ความดันโลหิต และน้ำหนักตัว พร้อมกับผลการวินิจฉัยว่าเป็นเบาหวานหรือไม่ จากนั้นตัวแบบจะเรียนรู้ความสัมพันธ์ระหว่างข้อมูลเหล่านี้กับผลลัพธ์ ในการเรียนรู้แบบไม่มีผู้สอน ตัวแบบจะได้รับเพียงข้อมูลนำเข้า $\mathbf{x}$ โดยไม่มีค่า $y$ ใด ๆ บอกกล่าว ดังนั้นตัวแบบจึงต้องค้นหาโครงสร้างและรูปแบบที่ซ่อนอยู่ในข้อมูลด้วยตนเอง

\begin{myLimit}[]{ความแตกต่างระหว่างการเรียนรู้แบบมีผู้สอนและไม่มีผู้สอน}{}
\begin{itemize}
    \item \textbf{การเรียนรู้แบบมีผู้สอน (Supervised Learning)} ต้องการข้อมูลที่มีการติดป้ายกำกับ (Labeled Data) ซึ่งมักต้องใช้แรงงานมนุษย์ในการสร้าง เป้าหมายคือการเรียนรู้การจับคู่ระหว่าง $\mathbf{x} \to y$ และสามารถวัดความผิดพลาดได้โดยตรงเมื่อเทียบกับคำตอบที่ถูกต้อง

    \item \textbf{การเรียนรู้แบบไม่มีผู้สอน (Unsupervised Learning)} ทำงานกับข้อมูลที่ไม่มีการติดป้ายกำกับ (Unlabeled Data) ซึ่งหาได้ง่ายและมีราคาถูกกว่ามาก เป้าหมายคือการค้นหาโครงสร้าง กลุ่ม หรือรูปแบบที่ซ่อนอยู่ในข้อมูล โดยไม่มี ``คำตอบที่ถูกต้อง'' เป็นตัวอ้างอิง
\end{itemize}
\end{myLimit}

ความท้าทายหลักของการเรียนรู้แบบไม่มีผู้สอนคือการประเมินประสิทธิภาพของตัวแบบ เนื่องจากไม่มีค่าความจริงพื้นฐาน (Ground Truth) ให้เปรียบเทียบ ในทางปฏิบัติ นักวิทยาศาสตร์ข้อมูลมักใช้เมทริกซ์ความคล้ายคลึง (Similarity Matrix) หรือค่าที่ได้จากฟังก์ชันวัตถุประสงค์ (Objective Function) เช่น ค่า Sum of Squared Errors (SSE) ในการจัดกลุ่ม หรือค่า Explained Variance ในการลดมิติ เพื่อประเมินและเปรียบเทียบผลลัพธ์ของตัวแบบ

งานหลักของการเรียนรู้แบบไม่มีผู้สอนสามารถจำแนกได้เป็น 3 ประเภทใหญ่ ๆ ดังนี้ ประเภทแรกคือ \textbf{การจัดกลุ่ม (Clustering)} ซึ่งมีเป้าหมายเพื่อแบ่งข้อมูลออกเป็นกลุ่มย่อย ๆ ตามความคล้ายคลึงกัน กลุ่มที่มีข้อมูลคล้ายกันจะถูกจัดไว้ในกลุ่มเดียวกัน ในขณะที่ข้อมูลที่แตกต่างกันจะอยู่ในกลุ่มที่ต่างกัน ประเภทที่สองคือ \textbf{การลดมิติ (Dimensionality Reduction)} ซึ่งมีเป้าหมายเพื่อแปลงข้อมูลที่มีมิติสูงให้เป็นข้อมูลที่มีมิติต่ำลงโดยพยายามรักษาข้อมูลสำคัญไว้ให้ได้มากที่สุด ประเภทที่สามคือ \textbf{การตรวจจับความผิดปกติ (Anomaly Detection)} ซึ่งมีเป้าหมายเพื่อระบุข้อมูลที่มีพฤติกรรมแตกต่างจากข้อมูลส่วนใหญ่อย่างมีนัยสำคัญ ซึ่งอาจบ่งบอกถึงปัญหาหรือเหตุการณ์ที่ต้องให้ความสนใจ

%####################################
\section{การจัดกลุ่มข้อมูล (Clustering)}

การจัดกลุ่มข้อมูล (Clustering) เป็นเทคนิคพื้นฐานที่สุดในการเรียนรู้แบบไม่มีผู้สอน โดยมีหลักการพื้นฐานคือการจัดกลุ่มวัตถุ (Objects) หรือข้อมูล (Data Points) ที่มีความคล้ายคลึงกันเข้าไว้ด้วยกัน ในขณะที่ข้อมูลที่แตกต่างกันจะถูกจัดไว้ในกลุ่มที่ต่างกัน การจัดกลุ่มมีการประยุกต์ใช้ในหลากหลายสาขา เช่น การแบ่งกลุ่มลูกค้าตามพฤติกรรมการซื้อสินค้า (Customer Segmentation) การจัดกลุ่มเอกสารตามหัวข้อเนื้อหา (Document Clustering) การจัดกลุ่มภาพถ่ายตามความคล้ายคลึงของเนื้อหา (Image Clustering) และการวิเคราะห์ประชากรศาสตร์ (Demographic Clustering) เป็นต้น

ความท้าทายสำคัญในการจัดกลุ่มข้อมูลอยู่ที่การนิยาม ``ความคล้ายคลึง'' (Similarity) หรือ ``ระยะทาง'' (Distance) ระหว่างข้อมูล ซึ่งโดยทั่วไปจะใช้ระยะทางแบบยุคลิด (Euclidean Distance) เป็นมาตรวัดหลัก สำหรับข้อมูลที่มีมิติสูง อาจพิจารณาใช้ระยะทางแมนฮัตตัน (Manhattan Distance) หรือเมทริกซ์ความคล้ายคลึงโคไซน์ (Cosine Similarity) แทน การเลือกมาตรวัดที่เหมาะสมขึ้นอยู่กับลักษณะของข้อมูลและเป้าหมายของการจัดกลุ่ม

ในส่วนนี้จะศึกษาวิธีการจัดกลุ่มที่ได้รับความนิยมสองวิธีหลัก ได้แก่ K-Means Clustering ซึ่งเป็นวิธีการจัดกลุ่มแบบแบ่งพาร์ติชัน (Partitioning Method) ที่ต้องกำหนดจำนวนกลุ่มล่วงหน้า และ Hierarchical Clustering ซึ่งเป็นวิธีการจัดกลุ่มแบบลำดับชั้น (Hierarchical Method) ที่สร้างโครงสร้างต้นไม้ (Dendrogram) ของการจัดกลุ่มโดยไม่ต้องกำหนดจำนวนกลุ่มล่วงหน้า

\subsection{K-Means Clustering}

K-Means Clustering เป็นอัลกอริทึมการจัดกลุ่มที่ได้รับความนิยมมากที่สุดในปัจจุบัน เนื่องจากมีความเรียบง่าย ทำความเข้าใจได้ง่าย และสามารถประมวลผลข้อมูลขนาดใหญ่ได้อย่างมีประสิทธิภาพ อัลกอริทึมนี้มีจุดกำเนิดจากแนวคิดการแบ่งข้อมูลออกเป็น $K$ กลุ่ม โดยแต่ละกลุ่มจะมีจุดศูนย์กลาง (Centroid) ที่เป็นค่าเฉลี่ยของข้อมูลในกลุ่มนั้น วัตถุประสงค์ของอัลกอริทึมคือการจัดข้อมูลแต่ละตัวไปยังกลุ่มที่ใกล้เคียงที่สุด เพื่อให้ค่าผลรวมของระยะทางยกกำลังสอง (Sum of Squared Errors หรือ SSE) มีค่าน้อยที่สุด

ก่อนเข้าสู่รายละเอียดของอัลกอริทึม K-Means จำเป็นต้องทำความเข้าใจพื้นฐานทางคณิตศาสตร์ที่เกี่ยวข้องกับระยะทางและการคำนวณจุดศูนย์กลางก่อน สำหรับข้อมูลที่มี $d$ มิติ ซึ่งแทนด้วยเวกเตอร์ $\mathbf{x}_i = (x_{i1}, x_{i2}, \ldots, x_{id})^\top \in \mathbb{R}^d$ ระยะทางแบบยุคลิด (Euclidean Distance) ระหว่างข้อมูลสองตัว $\mathbf{x}_i$ และ $\mathbf{x}_j$ นิยามโดย

\begin{equation}
d(\mathbf{x}_i, \mathbf{x}_j) = \|\mathbf{x}_i - \mathbf{x}_j\| = \sqrt{\sum_{k=1}^{d} (x_{ik} - x_{jk})^2}
\end{equation}

เมื่อพิจารณากลุ่ม $C$ ที่มีข้อมูล $n_c$ ตัว จุดศูนย์กลางของกลุ่ม (Centroid) $\boldsymbol{\mu}_c$ คำนวณได้จากค่าเฉลี่ยเวกเตอร์ของข้อมูลทั้งหมดในกลุ่ม

\begin{equation}
\boldsymbol{\mu}_c = \frac{1}{n_c} \sum_{\mathbf{x}_i \in C_c} \mathbf{x}_i
\end{equation}

โดยที่ $C_c$ แทนเซตของข้อมูลที่อยู่ในกลุ่ม $c$ และ $n_c = |C_c|$ คือจำนวนข้อมูลในกลุ่ม $c$ ค่า SSE (Sum of Squared Errors) ของกลุ่ม $c$ คือผลรวมของระยะทางยกกำลังสองระหว่างข้อมูลแต่ละตัวในกลุ่มกับจุดศูนย์กลางของกลุ่ม

\begin{equation}
\text{SSE}_c = \sum_{\mathbf{x}_i \in C_c} \|\mathbf{x}_i - \boldsymbol{\mu}_c\|^2
\end{equation}

และค่า SSE รวมของทุกกลุ่มคือ $\text{SSE} = \sum_{c=1}^{K} \text{SSE}_c$ ซึ่งเป็นค่าที่อัลกอริทึม K-Means พยายามทำให้มีค่าน้อยที่สุด

\begin{myBox}[]{ตัวแบบและฟังก์ชันวัตถุประสงค์ของ K-Means}{}
สมมติว่ามีข้อมูล $N$ ตัว แทนด้วยเซต $\mathcal{D} = \{\mathbf{x}_1, \mathbf{x}_2, \ldots, \mathbf{x}_N\}$ โดยแต่ละ $\mathbf{x}_i \in \mathbb{R}^d$ และต้องการแบ่งข้อมูลออกเป็น $K$ กลุ่ม ตัวแบบ K-Means จะกำหนดการจัดสรร (Assignment) $\mathbf{z}_i \in \{1, 2, \ldots, K\}$ สำหรับแต่ละข้อมูล $\mathbf{x}_i$ โดยที่ $\mathbf{z}_i = c$ หมายความว่า $\mathbf{x}_i$ ถูกจัดไปอยู่ในกลุ่ม $c$

ฟังก์ชันวัตถุประสงค์ (Objective Function) ที่ต้องการทำให้มีค่าน้อยที่สุดคือ

$$J = \sum_{c=1}^{K} \sum_{\mathbf{x}_i \in C_c} \|\mathbf{x}_i - \boldsymbol{\mu}_c\|^2$$

\noindent โดยที่ $\boldsymbol{\mu}_c$ คือจุดศูนย์กลางของกลุ่ม $c$
\end{myBox}

อัลกอริทึม K-Means ทำงานผ่านกระบวนการทำซ้ำสองขั้นตอนหลักที่เรียกว่า \textbf{Expectation-Maximization (EM)} จนกระทั่งค่าจุดศูนย์กลางของกลุ่มเปลี่ยนแปลงน้อยมากหรือจำนวนรอบถึงค่าสูงสุดที่กำหนด ขั้นตอนที่หนึ่งเรียกว่า \textbf{E-step} (Expectation Step) เป็นการจัดสรรข้อมูลแต่ละตัวไปยังกลุ่มที่มีจุดศูนย์กลางใกล้เคียงที่สุด ขั้นตอนที่สองเรียกว่า \textbf{M-step} (Maximization Step) เป็นการคำนวณจุดศูนย์กลางใหม่ของแต่ละกลุ่มจากข้อมูลที่ถูกจัดสรรเข้ามาในขั้นตอน E-step

\begin{exmp}
เพื่อให้เข้าใจการทำงานของอัลกอริทึม K-Means อย่างชัดเจน พิจารณาตัวอย่างการจัดกลุ่มข้อมูลลูกค้า 2 มิติ ซึ่งประกอบด้วยข้อมูล 6 รายการดังแสดงในตาราง \ref{tab:kmeans_data} โดยแต่ละลูกค้ามีคะแนนความภักดี (Loyalty Score) และยอดซื้อรายเดือน (Monthly Purchase) เป็นตัวแปร

\begin{table}[h!]
    \centering
    \caption{ข้อมูลตัวอย่างลูกค้า 6 รายสำหรับการจัดกลุ่มด้วย K-Means}
    \label{tab:kmeans_data}
    \begin{tabular}{|c|c|c|}
        \hline
        \textbf{ลูกค้า} & \textbf{คะแนนภักดี (Loyalty)} & \textbf{ยอดซื้อ (บาท)} \\
        \hline
        1 & 2 & 3 \\
        \hline
        2 & 4 & 7 \\
        \hline
        3 & 5 & 5 \\
        \hline
        4 & 8 & 8 \\
        \hline
        5 & 9 & 6 \\
        \hline
        6 & 7 & 4 \\
        \hline
    \end{tabular}
\end{table}

\textbf{ขั้นตอนที่ 1: กำหนดจำนวนกลุ่มและเลือกจุดศูนย์กลางเริ่มต้น}

สมมติกำหนด $K = 2$ กลุ่ม และสุ่มเลือกจุดศูนย์กลางเริ่มต้นเป็น $\boldsymbol{\mu}_1 = (2, 3)$ สำหรับกลุ่มที่ 1 และ $\boldsymbol{\mu}_2 = (8, 8)$ สำหรับกลุ่มที่ 2

\textbf{ขั้นตอนที่ 2: E-step --- จัดสรรข้อมูลไปยังกลุ่มที่ใกล้ที่สุด}

คำนวณระยะทางยุคลิดจากแต่ละลูกค้าไปยังจุดศูนย์กลางทั้งสอง แล้วจัดสรรไปยังกลุ่มที่มีระยะทางน้อยกว่า

สำหรับลูกค้าที่ 1: $d(\mathbf{x}_1, \boldsymbol{\mu}_1) = \sqrt{(2-2)^2 + (3-3)^2} = 0$ และ $d(\mathbf{x}_1, \boldsymbol{\mu}_2) = \sqrt{(2-8)^2 + (3-8)^2} = \sqrt{36+25} = \sqrt{61} \approx 7.81$ ดังนั้นลูกค้าที่ 1 อยู่ในกลุ่ม 1

สำหรับลูกค้าที่ 2: $d(\mathbf{x}_2, \boldsymbol{\mu}_1) = \sqrt{(4-2)^2 + (7-3)^2} = \sqrt{4+16} = \sqrt{20} \approx 4.47$ และ $d(\mathbf{x}_2, \boldsymbol{\mu}_2) = \sqrt{(4-8)^2 + (7-8)^2} = \sqrt{16+1} = \sqrt{17} \approx 4.12$ ดังนั้นลูกค้าที่ 2 อยู่ในกลุ่ม 2

ทำเช่นนี้กับทุกลูกค้า ผลลัพธ์คือ กลุ่ม 1 ประกอบด้วยลูกค้า $\{1, 3, 6\}$ และกลุ่ม 2 ประกอบด้วยลูกค้า $\{2, 4, 5\}$

\textbf{ขั้นตอนที่ 3: M-step --- คำนวณจุดศูนย์กลางใหม่}

จุดศูนย์กลางกลุ่ม 1: $\boldsymbol{\mu}_1 = \frac{1}{3}[(2,3) + (5,5) + (7,4)] = \left(\frac{14}{3}, \frac{12}{3}\right) = (4.67, 4)$

จุดศูนย์กลางกลุ่ม 2: $\boldsymbol{\mu}_2 = \frac{1}{3}[(4,7) + (8,8) + (9,6)] = \left(\frac{21}{3}, \frac{21}{3}\right) = (7, 7)$

\textbf{ขั้นตอนที่ 4: ทำซ้ำ E-step และ M-step}

E-step ใหม่: ลูกค้าที่ 2 มีระยะทางถึง $\boldsymbol{\mu}_1 = (4.67, 4)$ เท่ากับ $\sqrt{(4-4.67)^2 + (7-4)^2} = \sqrt{0.45 + 9} \approx 3.07$ และระยะทางถึง $\boldsymbol{\mu}_2 = (7, 7)$ เท่ากับ $\sqrt{(4-7)^2 + (7-7)^2} = 3$ ดังนั้นลูกค้าที่ 2 ยังคงอยู่ในกลุ่ม 2 เนื่องจากระยะทางถึง $\boldsymbol{\mu}_2$ น้อยกว่า

อัลกอริทึมจะหยุดเมื่อการจัดสรรไม่เปลี่ยนแปลงอีก หรือจำนวนรอบถึงค่าสูงสุด
\end{exmp}

การเลือกค่า $K$ ที่เหมาะสมเป็นปัญหาสำคัญในการใช้งาน K-Means วิธีการหนึ่งที่ได้รับความนิยมคือ \textbf{Elbow Method} ซึ่งพล็อตกราฟระหว่างค่า $K$ กับค่า SSE แล้วมองหาจุดที่กราฟเปลี่ยนจากความชันมากเป็นความชันน้อย ซึ่งมักมีลักษณะคล้ายข้อศอก (Elbow) ของแขน จุดนั้นบ่งบอกว่าการเพิ่มจำนวนกลุ่มมากขึ้นไม่ช่วยลดค่า SSE อย่างมีนัยสำคัญ

นอกจากนี้ ยังมีวิธี \textbf{Silhouette Score} ที่วัดว่าข้อมูลแต่ละตัว ``เข้ากับกลุ่มของตัวเองได้ดีเพียงใด'' เมื่อเทียบกับ ``กลุ่มข้างเคียงที่ใกล้ที่สุด'' ค่า Silhouette Score อยู่ระหว่าง $-1$ ถึง $+1$ โดยค่าที่ใกล้ $+1$ บ่งบอกว่าข้อมูลถูกจัดไปยังกลุ่มที่ถูกต้อง ค่าที่ใกล้ $0$ บ่งบอกว่าข้อมูลอยู่บนเขตแดนระหว่างกลุ่ม และค่าที่ใกล้ $-1$ บ่งบอกว่าข้อมูลถูกจัดไปยังกลุ่มที่ไม่ถูกต้อง

การประมวลผล K-Means ด้วยภาษา Python แสดงดังภาพ \ref{fig:kmeans_python} ซึ่งจะสังเกตได้ว่าการใช้ \texttt{sklearn.cluster.KMeans} ทำให้สามารถสร้างตัวแบบและทำนายการจัดกลุ่มได้อย่างสะดวก

\begin{figure}[h]
    \centering
    \includegraphics[width=0.9\textwidth]{kmeans_clustering_example.png}
    \caption{ผลลัพธ์การจัดกลุ่มด้วย K-Means พร้อมแสดงจุดศูนย์กลางของแต่ละกลุ่ม}
    \label{fig:kmeans_result}
\end{figure}

ข้อจำกัดสำคัญของ K-Means ที่ควรตระหนักมีหลายประการ ประการแรก K-Means ตั้งสมมติฐานว่ากลุ่มมีรูปร่างเป็นทรงกลม (Spherical) และขนาดใกล้เคียงกัน ซึ่งอาจไม่เป็นจริงในข้อมูลจริงบางชุด ประการที่สอง K-Means ไวต่อค่าเริ่มต้นของจุดศูนย์กลาง ซึ่งอาจทำให้ติดอยู่ในค่าเฉพาะท้องถิ่น (Local Optimum) ที่ไม่ใช่ค่าที่ดีที่สุด ประการที่สาม K-Means ไม่สามารถจัดการกับข้อมูลที่มีมิติสูงมาก ๆ ได้ดีนักเนื่องจากปัญหา ``มิติเสExecute_str ียน'' (Curse of Dimensionality) และประการสุดท้าย K-Means ไม่ deterministic ซึ่งหมายความว่าผลลัพธ์อาจแตกต่างกันในแต่ละครั้งที่รัน

\begin{Verbatim}[tabsize=4]
# K-Means Clustering in Python
from sklearn.cluster import KMeans
import numpy as np
import matplotlib.pyplot as plt

# Sample data: loyalty score and monthly purchase (in thousands)
X = np.array([[2, 3], [4, 7], [5, 5], [8, 8], [9, 6], [7, 4]])

# Apply K-Means with K=2
kmeans = KMeans(n_clusters=2, init='random', n_init=10, random_state=42)
labels = kmeans.fit_predict(X)
centers = kmeans.cluster_centers_

# Plot results
plt.scatter(X[:, 0], X[:, 1], c=labels, cmap='viridis', s=100)
plt.scatter(centers[:, 0], centers[:, 1], c='red', marker='X', s=200)
plt.xlabel('Loyalty Score')
plt.ylabel('Monthly Purchase (K THB)')
plt.title('K-Means Clustering Result (K=2)')
plt.show()

# Print SSE for different K values using Elbow Method
for k in range(1, 5):
    kmeans = KMeans(n_clusters=k, n_init=10, random_state=42)
    kmeans.fit(X)
    print(f"K={k}, SSE={kmeans.inertia_:.2f}")
\end{Verbatim}

\subsection{Hierarchical Clustering}

Hierarchical Clustering เป็นอีกหนึ่งวิธีการจัดกลุ่มที่ได้รับความนิยม โดยมีความแตกต่างจาก K-Means ตรงที่ไม่ต้องกำหนดจำนวนกลุ่ม $K$ ล่วงหน้า ผลลัพธ์ของ Hierarchical Clustering อยู่ในรูปของโครงสร้างต้นไม้ (Dendrogram) ที่แสดงความสัมพันธ์แบบลำดับชั้นระหว่างข้อมูล ทำให้สามารถสำรวจโครงสร้างของข้อมูลได้อย่างลึกซึ้งและเข้าใจความคล้ายคลึงระหว่างกลุ่มในหลายระดับ

Hierarchical Clustering แบ่งออกเป็น 2 ประเภทหลักตามทิศทางการสร้างโครงสร้าง ประเภทแรกคือ \textbf{Agglomerative (Bottom-Up)} ซึ่งเริ่มต้นด้วยการถือว่าแต่ละข้อมูลเป็นกลุ่มของตัวเอง (แต่ละจุดคือกลุ่ม) แล้วค่อย ๆ รวมกลุ่ม (Merge) ที่ใกล้เคียงกันเข้าด้วยกันจนกลายเป็นกลุ่มเดียว ประเภทที่สองคือ \textbf{Divisive (Top-Down)} ซึ่งเริ่มต้นด้วยกลุ่มเดียวที่ประกอบด้วยทุกข้อมูล แล้วค่อย ๆ แยก (Split) ออกเป็นกลุ่มย่อย ๆ จนกลายเป็นกลุ่มละหนึ่งจุด ในทางปฏิบัติ Agglomerative ถูกใช้มากกว่าเนื่องจากมีการนำไปใช้งานที่ง่ายกว่า

หัวใจสำคัญของ Agglomerative Clustering คือการนิยามระยะทางระหว่างกลุ่ม (Linkage Criteria) ซึ่งมีหลายวิธีดังนี้

\begin{myBox}[]{วิธีการคำนวณระยะทางระหว่างกลุ่ม (Linkage Methods)}{}
ให้ $C_i$ และ $C_j$ เป็นกลุ่มสองกลุ่ม และ $d(\mathbf{x}, \mathbf{y})$ เป็นระยะทางระหว่างข้อมูลสองตัว วิธีการคำนวณระยะทางระหว่างกลุ่มมีดังนี้

\begin{enumerate}
    \item \textbf{Single Linkage:} ระยะทางระหว่างกลุ่มคือระยะทางที่น้อยที่สุดระหว่างจุดในกลุ่ม $C_i$ กับจุดในกลุ่ม $C_j$
    $$d_{\text{single}}(C_i, C_j) = \min_{\mathbf{x} \in C_i, \mathbf{y} \in C_j} d(\mathbf{x}, \mathbf{y})$$

    \item \textbf{Complete Linkage:} ระยะทางระหว่างกลุ่มคือระยะทางที่มากที่สุดระหว่างจุดในกลุ่ม $C_i$ กับจุดในกลุ่ม $C_j$
    $$d_{\text{complete}}(C_i, C_j) = \max_{\mathbf{x} \in C_i, \mathbf{y} \in C_j} d(\mathbf{x}, \mathbf{y})$$

    \item \textbf{Average Linkage:} ระยะทางระหว่างกลุ่มคือค่าเฉลี่ยของระยะทางระหว่างทุกคู่ของจุดในกลุ่ม $C_i$ กับ $C_j$
    $$d_{\text{average}}(C_i, C_j) = \frac{1}{|C_i| \cdot |C_j|} \sum_{\mathbf{x} \in C_i} \sum_{\mathbf{y} \in C_j} d(\mathbf{x}, \mathbf{y})$$

    \item \textbf{Ward's Method:} เพิ่มค่า Sum of Squared Errors (SSE) ที่เกิดจากการรวมกลุ่ม รวมกลุ่มที่ทำให้ SSE เพิ่มขึ้นน้อยที่สุด
\end{enumerate}
\end{myBox}

แต่ละวิธีการมีคุณสมบัติและข้อจำกัดที่แตกต่างกัน Single Linkage มีแนวโน้มที่จะสร้างกลุ่มที่ยาวเรียว (Strand-like Clusters) แต่ทำงานได้ดีกับข้อมูลที่มีรูปร่างซับซ้อน Complete Linkage มักสร้างกลุ่มที่มีขนาดกลมเกลี้ยงกลมเกือบกัน แต่ไวต่อ Outliers Average Linkage เป็นวิธีกลาง ๆ ที่สมดุลระหว่างความเรียบง่ายและความทนทาน Ward's Method ให้ผลลัพธ์ที่ดีเมื่อกลุ่มมีขนาดใกล้เคียงกันและมีรูปร่างทรงกลม

\begin{exmp}
เพื่อทำความเข้าใจ Hierarchical Clustering อย่างเป็นรูปธรรม พิจารณาข้อมูลลูกค้า 4 รายดังนี้: $\mathbf{x}_1 = (1, 1)$, $\mathbf{x}_2 = (1.5, 2)$, $\mathbf{x}_3 = (5, 6)$, $\mathbf{x}_4 = (6, 5)$ จะทำการจัดกลุ่มด้วย Agglomerative Hierarchical Clustering โดยใช้ Single Linkage

\textbf{ขั้นตอนที่ 1: เริ่มต้น} แต่ละจุดเป็นกลุ่มของตัวเอง: $\{1\}, \{2\}, \{3\}, \{4\}$

\textbf{ขั้นตอนที่ 2: รอบที่ 1 --- หาคู่ที่ใกล้ที่สุด}

คำนวณระยะทาง Single Linkage:
\begin{align*}
d(\{1\}, \{2\}) &= \sqrt{(1-1.5)^2 + (1-2)^2} = \sqrt{0.25 + 1} = \sqrt{1.25} \approx 1.12 \\
d(\{1\}, \{3\}) &= \sqrt{(1-5)^2 + (1-6)^2} = \sqrt{16 + 25} = \sqrt{41} \approx 6.40 \\
d(\{1\}, \{4\}) &= \sqrt{(1-6)^2 + (1-5)^2} = \sqrt{25 + 16} = \sqrt{41} \approx 6.40 \\
d(\{2\}, \{3\}) &= \sqrt{(1.5-5)^2 + (2-6)^2} = \sqrt{12.25 + 16} = \sqrt{28.25} \approx 5.32 \\
d(\{2\}, \{4\}) &= \sqrt{(1.5-6)^2 + (2-5)^2} = \sqrt{20.25 + 9} = \sqrt{29.25} \approx 5.41 \\
d(\{3\}, \{4\}) &= \sqrt{(5-6)^2 + (6-5)^2} = \sqrt{1 + 1} = \sqrt{2} \approx 1.41
\end{align*}

คู่ที่ใกล้ที่สุดคือ $\{1\}$ กับ $\{2\}$ ด้วยระยะทาง $1.12$ รวมเป็นกลุ่ม $\{1, 2\}$

\textbf{ขั้นตอนที่ 3: รอบที่ 2 --- คำนวณใหม่ด้วย Single Linkage}

ระยะทางจากกลุ่ม $\{1, 2\}$ ไปยัง $\{3\}$ คือ $\min(d(\{1\}, \{3\}), d(\{2\}, \{3\})) = \min(6.40, 5.32) = 5.32$

ระยะทางจากกลุ่ม $\{1, 2\}$ ไปยัง $\{4\}$ คือ $\min(d(\{1\}, \{4\}), d(\{2\}, \{4\})) = \min(6.40, 5.41) = 5.41$

ระยะทางระหว่าง $\{3\}$ กับ $\{4\}$ คือ $1.41$ ซึ่งเป็นค่าน้อยที่สุด ดังนั้นรวม $\{3\}$ กับ $\{4\}$ เป็นกลุ่ม $\{3, 4\}$

\textbf{ขั้นตอนที่ 4: รอบสุดท้าย --- รวมกลุ่มใหญ่}

ระยะทางระหว่าง $\{1, 2\}$ กับ $\{3, 4\}$ คือ $\min(5.32, 5.41) = 5.32$ รวมเป็นกลุ่มเดียว $\{1, 2, 3, 4\}$

ผลลัพธ์คือโครงสร้างลำดับชั้นที่แสดงด้วย Dendrogram
\end{exmp}

\begin{figure}[h]
    \centering
    \includegraphics[width=0.8\textwidth]{hierarchical_clustering_dendrogram.png}
    \caption{Dendrogram แสดงโครงสร้างลำดับชั้นของการจัดกลุ่มข้อมูลลูกค้า 4 ราย}
    \label{fig:dendrogram}
\end{figure}

ข้อดีสำคัญของ Hierarchical Clustering คือสามารถสร้าง Dendrogram ที่แสดงโครงสร้างลำดับชั้นของข้อมูล ทำให้ผู้วิเคราะห์สามารถสำรวจความสัมพันธ์ระหว่างข้อมูลในหลายระดับของความลึกได้ อย่างไรก็ตาม ข้อเสียคือความซับซ้อนในการคำนวณ เนื่องจากในกรณีที่แย่ที่สุด อัลกอริทึมมีความซับซ้อน $O(n^2 \log n)$ ซึ่งทำให้ไม่เหมาะกับข้อมูลขนาดใหญ่มาก ๆ เมื่อเทียบกับ K-Means ที่มีความซับซ้อนเพียง $O(nK)$ ต่อรอบ

%####################################
\section{การลดมิติข้อมูล (Dimensionality Reduction)}

ข้อมูลในโลกจริงมักมีมิติ (Dimensions) จำนวนมาก ข้อมูลลูกค้าอาจมีหลายร้อยคุณลักษณะ ข้อมูลภาพอาจมีหลายพันพิกเซล และข้อมูลข้อความอาจมีหลายหมื่นคำศัพท์ ปัญหา ``มิติเสExecute_strียน'' (Curse of Dimensionality) บ่งชี้ว่าเมื่อจำนวนมิติเพิ่มขึ้น ข้อมูลจะกระจายตัวมากขึ้นในปริภูมิ ทำให้การค้นหารูปแบบและการเรียนรู้ตัวแบบกลายเป็นเรื่องยากและต้องการข้อมูลจำนวนมากขึ้นแบบทวีคูณ การลดมิติข้อมูล (Dimensionality Reduction) จึงเป็นเทคนิคสำคัญที่ช่วยแก้ปัญหานี้โดยการแปลงข้อมูลจากปริภูมิมิติสูงไปยังปริภูมิมิติต่ำกว่า โดยพยายามรักษาข้อมูลสำคัญและโครงสร้างของข้อมูลไว้ให้ได้มากที่สุด

การลดมิติข้อมูลสามารถจำแนกออกเป็นสองประเภทหลัก ประเภทแรกคือ \textbf{การเลือกคุณลักษณะ (Feature Selection)} ซึ่งเป็นการเลือกเฉพาะคุณลักษณะบางตัวจากชุดคุณลักษณะเดิม โดยไม่ทำการแปลงข้อมูล ตัวอย่างเช่น การใช้ Correlation Matrix เพื่อคัดเลือกคุณลักษณะที่มีความสัมพันธ์กับตัวแปรเป้าหมายสูง หรือการใช้ Recursive Feature Elimination (RFE) ประเภทที่สองคือ \textbf{การสกัดคุณลักษณะ (Feature Extraction)} ซึ่งเป็นการสร้างคุณลักษณะใหม่จากคุณลักษณะเดิมโดยการรวมหรือแปลงข้อมูล ตัวอย่างเช่น Principal Component Analysis (PCA) ที่จะกล่าวถึงในรายละเอียด

การลดมิติข้อมูลมีประโยชน์หลายประการ ประการแรกช่วยลดความซับซ้อนในการคำนวณและเวลาในการฝึกตัวแบบ ประการที่สองช่วยลดปัญหา Overfitting โดยการลด Noise ในข้อมูล ประการที่สามช่วยให้สามารถแสดงผลข้อมูลในมิติที่มนุษย์เข้าใจได้ (เช่น 2 มิติ หรือ 3 มิติ) และประการสุดท้ายช่วยประหยัดพื้นที่ในการจัดเก็บข้อมูล

\subsection{Principal Component Analysis (PCA)}

Principal Component Analysis (PCA) เป็นเทคนิคการลดมิติข้อมูลที่ได้รับความนิยมมากที่สุดในปัจจุบัน PCA ทำงานโดยการค้นหาแกน (Axes) ใหม่ในปริภูมิมิติสูงที่เรียกว่า \textbf{Principal Components (PCs)} ซึ่งแต่ละแกนจะถูกจัดเรียงตามความสำคัญ (Variance) ที่อธิบายได้ แกนแรก (PC1) จะอธิบายความแปรปรวนของข้อมูลได้มากที่สุด แกนที่สอง (PC2) จะอธิบายความแปรปรวนที่เหลือได้มากที่สุดหลังจาก PC1 และเป็นไปในทิศทางที่ตั้งฉาก (Orthogonal) กับ PC1 ทำให้แกนใหม่ทั้งหมดตั้งฉากซึ่งกันและกัน

หลักการพื้นฐานของ PCA สามารถอธิบายได้ด้วยแนวคิดทางเรขาคณิต สมมติมีข้อมูล 2 มิติที่มีการกระจายตัวเป็นรูปวงรี PCA จะหาแกนหลัก (Major Axis) ของวงรีนั้น ซึ่งเป็นทิศทางที่ข้อมูลมีความแปรปรวนมากที่สุด แกนนี้จะกลายเป็น PC1 จากนั้น PCA จะหาแกนรอง (Minor Axis) ที่ตั้งฉากกับ PC1 ซึ่งจะเป็น PC2 เมื่อฉาย (Project) ข้อมูลลงบน PC1 และ PC2 จะได้ข้อมูลใหม่ที่มี 2 มิติเช่นเดิม แต่หากต้องการลดเหลือ 1 มิติ ก็สามารถฉายลงบน PC1 เพียงแกนเดียว

\begin{myBox}[]{นิยามทางคณิตศาสตร์ของ PCA}{}
ให้ $\mathbf{X} \in \mathbb{R}^{N \times d}$ เป็นเมทริกซ์ข้อมูล โดยที่ $N$ คือจำนวนตัวอย่าง และ $d$ คือจำนวนมิติ (คุณลักษณะ) ขั้นตอนการทำ PCA มีดังนี้

\textbf{ขั้นตอนที่ 1: ทำให้เป็นศูนย์กลาง (Centering)} คำนวณค่าเฉลี่ยของแต่ละคุณลักษณะ แล้วลบออกจากข้อมูล
$$\mathbf{x}_i' = \mathbf{x}_i - \bar{\mathbf{x}}, \quad \bar{\mathbf{x}} = \frac{1}{N} \sum_{i=1}^{N} \mathbf{x}_i$$

\textbf{ขั้นตอนที่ 2: คำนวณเมทริกซ์ความแปรปรวนร่วม (Covariance Matrix)}
$$\mathbf{S} = \frac{1}{N-1} \mathbf{X}'^\top \mathbf{X}'$$

\textbf{ขั้นตอนที่ 3: หา Eigenvectors และ Eigenvalues} ของเมทริกซ์ความแปรปรวนร่วม $\mathbf{S}$ ซึ่งเป็น Principal Components

$$\mathbf{S} \mathbf{v}_j = \lambda_j \mathbf{v}_j$$

โดยที่ $\lambda_1 \geq \lambda_2 \geq \cdots \geq \lambda_d \geq 0$ คือ Eigenvalues และ $\mathbf{v}_j$ คือ Eigenvectors ที่ตรงกัน

\textbf{ขั้นตอนที่ 4: เลือก $k$ Principal Components} ที่มี Eigenvalues สูงสุด แล้วฉายข้อมูลลงบนแกนเหล่านั้น

$$\mathbf{z}_i = \mathbf{V}_k^\top \mathbf{x}_i'$$

โดยที่ $\mathbf{V}_k = [\mathbf{v}_1, \mathbf{v}_2, \ldots, \mathbf{v}_k] \in \mathbb{R}^{d \times k}$ และ $\mathbf{z}_i \in \mathbb{R}^k$ คือข้อมูลที่ถูกลดมิติแล้ว
\end{myBox}

คุณสมบัติสำคัญประการหนึ่งของ PCA คือ Eigenvalue $\lambda_j$ ที่ตรงกับ Principal Component ที่ $j$ จะเท่ากับค่าความแปรปรวนของข้อมูลที่ถูกฉายลงบนแกนนั้น ดังนั้นสัดส่วนของความแปรปรวนที่ PC $j$ อธิบายได้คือ $\lambda_j / \sum_{k=1}^{d} \lambda_k$ และสัดส่วนของความแปรปรวนรวมที่ $k$ Principal Components แรกอธิบายได้คือ $\sum_{j=1}^{k} \lambda_j / \sum_{j=1}^{d} \lambda_j$ ซึ่งเป็นค่าที่ใช้ในการตัดสินใจว่าควรรักษา Principal Components กี่ตัว

\begin{exmp}
เพื่อทำความเข้าใจ PCA อย่างเป็นรูปธรรม พิจารณาข้อมูล 2 มิติที่มี 5 ตัวอย่าง: $\mathbf{x}_1 = (2, 3)$, $\mathbf{x}_2 = (3, 6)$, $\mathbf{x}_3 = (4, 5)$, $\mathbf{x}_4 = (5, 8)$, $\mathbf{x}_5 = (6, 7)$ จะทำ PCA เพื่อลดมิติจาก 2 มิติเป็น 1 มิติ

\textbf{ขั้นตอนที่ 1: Centering}

ค่าเฉลี่ย: $\bar{x}_1 = (2+3+4+5+6)/5 = 4$ และ $\bar{x}_2 = (3+6+5+8+7)/5 = 5.8$

ข้อมูลที่ทำให้เป็นศูนย์กลาง:
\begin{align*}
\mathbf{x}_1' &= (-2, -2.8) \\
\mathbf{x}_2' &= (-1, 0.2) \\
\mathbf{x}_3' &= (0, -0.8) \\
\mathbf{x}_4' &= (1, 2.2) \\
\mathbf{x}_5' &= (2, 1.2)
\end{align*}

\textbf{ขั้นตอนที่ 2: คำนวณ Covariance Matrix}

เนื่องจากข้อมูล 2 มิติ:
$$\mathbf{S} = \begin{pmatrix} s_{11} & s_{12} \\ s_{21} & s_{22} \end{pmatrix}$$

โดยที่ $s_{11} = \frac{1}{4}[(-2)^2 + (-1)^2 + 0^2 + 1^2 + 2^2] = \frac{10}{4} = 2.5$

$s_{22} = \frac{1}{4}[(-2.8)^2 + (0.2)^2 + (-0.8)^2 + (2.2)^2 + (1.2)^2] = \frac{15.6}{4} = 3.9$

$s_{12} = s_{21} = \frac{1}{4}[(-2)(-2.8) + (-1)(0.2) + (0)(-0.8) + (1)(2.2) + (2)(1.2)] = \frac{10.2}{4} = 2.55$

ดังนั้น $\mathbf{S} = \begin{pmatrix} 2.5 & 2.55 \\ 2.55 & 3.9 \end{pmatrix}$

\textbf{ขั้นตอนที่ 3: หา Eigenvalues}

แก้สมการ $\det(\mathbf{S} - \lambda\mathbf{I}) = 0$:
$$\det\begin{pmatrix} 2.5-\lambda & 2.55 \\ 2.55 & 3.9-\lambda \end{pmatrix} = (2.5-\lambda)(3.9-\lambda) - 2.55^2 = 0$$

$$\lambda^2 - 6.4\lambda + 2.8875 = 0$$

$$\lambda = \frac{6.4 \pm \sqrt{6.4^2 - 4 \times 2.8875}}{2} = \frac{6.4 \pm \sqrt{23.425}}{2} = \frac{6.4 \pm 4.84}{2}$$

ได้ $\lambda_1 = 5.62$ และ $\lambda_2 = 0.78$ ซึ่งบ่งบอกว่า PC1 อธิบายความแปรปรวนได้ $5.62/(5.62+0.78) = 87.8\%$ และ PC2 อธิบายได้ $12.2\%$

การลดมิติจาก 2 เป็น 1 โดยใช้เฉพาะ PC1 จะสูญเสียข้อมูลเพียง $12.2\%$ เท่านั้น ซึ่งถือว่าคุ้มค่า
\end{exmp}

PCA มีการประยุกต์ใช้งานหลากหลาย ในด้านการวิเคราะห์ภาพ PCA สามารถใช้ในการบีบอัดภาพ (Image Compression) โดยการลดจำนวน Principal Components จากหลายร้อยพิกเซลเหลือเพียงไม่กี่สิบค่า ซึ่งยังคงรักษาลักษณะสำคัญของภาพไว้ได้ดี ในด้านการวิเคราะห์ข้อมูลทางการเงิน PCA สามารถใช้ในการสร้างดัชนีร่วม (Composite Index) จากตัวชี้วัดหลายตัวที่มีความสัมพันธ์กัน เช่น การสร้างดัชนีตลาดจากราคาหุ้นหลายตัว ในด้านการจำแนกประเภท PCA สามารถใช้เป็นขั้นตอนเตรียมข้อมูล (Preprocessing) ก่อนนำไปใช้กับตัวแบบ Classification เพื่อลด Noise และความซับซ้อนของข้อมูล

\begin{Verbatim}[tabsize=4]
# PCA in Python
from sklearn.decomposition import PCA
from sklearn.preprocessing import StandardScaler
import numpy as np

# Sample data: 5 customers with 2 features (loyalty and purchase)
X = np.array([[2, 3], [3, 6], [4, 5], [5, 8], [6, 7]])

# Step 1: Standardize the data (important for PCA)
scaler = StandardScaler()
X_scaled = scaler.fit_transform(X)

# Step 2: Apply PCA
pca = PCA(n_components=1)  # Reduce to 1 dimension
X_pca = pca.fit_transform(X_scaled)

print(f"Explained variance ratio: {pca.explained_variance_ratio_[0]:.4f}")
print(f"Reduced data:\n{X_pca}")
print(f"Original shape: {X.shape} -> Reduced shape: {X_pca.shape}")

# Check explained variance for different components
pca_full = PCA().fit(X_scaled)
print(f"Cumulative explained variance: {np.cumsum(pca_full.explained_variance_ratio_)}")
\end{Verbatim}

ข้อจำกัดสำคัญของ PCA มีหลายประการ ประการแรก PCA ตั้งสมมติฐานว่าความสัมพันธ์ระหว่างตัวแปรเป็นแบบเชิงเส้น (Linear) หากข้อมูลมีความสัมพันธ์แบบไม่เชิงเส้น (Non-linear) เช่น ข้อมูลที่มีโครงสร้างเป็นวงกลมหรือเกลียว PCA อาจไม่สามารถลดมิติได้อย่างมีประสิทธิภาพ ประการที่สอง PCA อ่อนไหวต่อ Outliers เนื่องจากการคำนวณ Covariance Matrix จะได้รับผลกระทบจากค่าที่ผิดปกติ ประการที่สาม PCA เป็น Unsupervised Method ซึ่งหมายความว่าไม่ได้พิจารณาตัวแปรเป้าหมาย (Label) ในการเลือก Principal Components ทำให้อาจลดมิติในทิศทางที่ไม่เกี่ยวข้องกับงานที่ต้องการ

สำหรับข้อมูลที่มีความสัมพันธ์แบบไม่เชิงเส้น มีเทคนิคการลดมิติอื่น ๆ ที่เหมาะสมกว่า เช่น t-SNE (t-Distributed Stochastic Neighbor Embedding) ซึ่งเป็นเทคนิคที่รักษาโครงสร้างท้องถิ่น (Local Structure) ของข้อมูลได้ดี หรือ UMAP (Uniform Manifold Approximation and Projection) ซึ่งทำงานได้เร็วกว่าและรักษาทั้งโครงสร้างท้องถิ่นและโครงสร้างระดับโลก (Global Structure) ได้ดีกว่า t-SNE

\subsection{การประยุกต์ใช้}

การลดมิติข้อมูลมีการประยุกต์ใช้งานอย่างกว้างขวางในหลากหลายสาขา ในด้านการประมวลผลภาพ (Image Processing) เทคนิค PCA ถูกนำไปใช้ในระบบจดจำใบหน้า (Face Recognition) โดยผ่านชื่อ Eigenfaces ซึ่งเป็น Principal Components ของภาพใบหน้าจำนวนมาก วิธีนี้ช่วยลดจำนวนมิติจากหลายพันพิกเซลเหลือเพียงไม่กี่ร้อยมิติโดยยังคงรักษาข้อมูลสำคัญในการจดจำใบหน้าไว้ได้ ในด้านการประมวลผลข้อความ (Text Processing) เทคนิคที่คล้ายกันถูกนำไปใช้กับเมทริกซ์คำศัพท์ (Term-Document Matrix) เพื่อลดมิติจากหลายหมื่นคำเหลือเพียงไม่กี่ร้อยมิติ ทำให้สามารถค้นหาเอกสารที่คล้ายคลึงกันได้อย่างมีประสิทธิภาพ

ในด้านการวิเคราะห์ข้อมูลจีโนม (Genomic Data Analysis) PCA ถูกนำไปใช้อย่างแพร่หลายในการตรวจสอบโครงสร้างประชากร (Population Structure) และการระบุ Outliers ทางพันธุกรรม โดยข้อมูลจีโนมของบุคคลจำนวนมากสามารถลดมิติจากหลายพัน SNPs (Single Nucleotide Polymorphisms) เหลือเพียง 2-3 มิติ เพื่อแสดงผลเป็นภาพที่เข้าใจง่าย ในด้านการเงิน PCA ถูกใช้ในการลดจำนวนปัจจัยเสี่ยง (Risk Factors) ที่มีความสัมพันธ์กัน เพื่อสร้างตัวแบบการบริหารความเสี่ยงที่เรียบง่ายขึ้น เช่น การสร้างดัชนีผลตอบแทนตลาดจากหุ้นจำนวนมาก

%####################################
\section{การตรวจจับความผิดปกติ (Anomaly Detection)}

การตรวจจับความผิดปกติ (Anomaly Detection) เป็นเทคนิคที่มีเป้าหมายเพื่อระบุข้อมูลที่มีพฤติกรรมแตกต่างอย่างมีนัยสำคัญจากข้อมูลส่วนใหญ่ ความผิดปกติ (Anomaly หรือ Outlier) คือข้อมูลที่ไม่เข้ากับรูปแบบปกติ (Normal Pattern) ของข้อมูลส่วนใหญ่ ซึ่งอาจบ่งบอกถึงปัญหาในระบบ การทุจริต การบุกรุก หรือเหตุการณ์ที่ต้องให้ความสนใจเป็นพิเศษ ตัวอย่างเช่น ธุรกรรมทางการเงินที่ผิดปกติอาจบ่งบอกถึงการฉ้อโกง สัญญาณเซ็นเซอร์ในโรงงานที่ผิดปกติอาจบ่งบอกถึงความเสี่ยงต่อการขัดข้องของเครื่องจักร และรูปแบบการเข้าใช้งานระบบคอมพิวเตอร์ที่ผิดปกติอาจบ่งบอกถึงการบุกรุกจากผู้ไม่หวังดี

ความแตกต่างระหว่างคำว่า Outlier กับ Anomaly มีความละเอียดอ่อน โดยทั่วไป Outlier หมายถึงข้อมูลที่มีค่าผิดปกติเมื่อเทียบกับการแจกแจงทางสถิติของข้อมูล ในขณะที่ Anomaly อาจหมายถึงรูปแบบ (Pattern) ที่ผิดปกติซึ่งอาจไม่จำเป็นต้องมีค่า экстремальный แต่มีลักษณะเฉพาะที่แตกต่างจากข้อมูลปกติ อย่างไรก็ตามในทางปฏิบัติสองคำนี้มักถูกใช้แทนกัน

วิธีการตรวจจับความผิดปกติสามารถจำแนกได้ตามแนวทางการเรียนรู้ดังนี้ วิธีแรกคือ \textbf{Statistical Methods} ซึ่งใช้การแจกแจงทางสถิติ (Statistical Distribution) ของข้อมูลปกติเป็นเกณฑ์ในการตัดสินว่าข้อมูลใหม่เป็นความผิดปกติหรือไม่ ตัวอย่างเช่น Z-Score ที่วัดว่าค่าห่างจากค่าเฉลี่ยกี่เท่าของส่วนเบี่ยงเบนมาตรฐาน หรือ IQR (Interquartile Range) ที่พิจารณาค่าที่อยู่นอกช่วง $Q_1 - 1.5 \times IQR$ ถึง $Q_3 + 1.5 \times IQR$ วิธีที่สองคือ \textbf{Proximity-Based Methods} ซึ่งใช้ระยะทางหรือความหนาแน่นในปริภูมิข้อมูลเป็นเกณฑ์ ข้อมูลที่อยู่ห่างจากข้อมูลอื่น ๆ มาก ๆ หรืออยู่ในบริเวณที่มีความหนาแน่นต่ำจะถูกระบุว่าเป็นความผิดปกติ ตัวอย่างเช่น K-Nearest Neighbors Distance หรือ Local Outlier Factor (LOF) วิธีที่สามคือ \textbf{Reconstruction-Based Methods} ซึ่งเรียนรู้ตัวแบบที่สามารถสร้าง (Reconstruct) ข้อมูลปกติได้ดี แล้วใช้ค่าความผิดพลาดในการสร้างใหม่เป็นเกณฑ์ในการตัดสิน ตัวอย่างเช่น Autoencoder ที่เรียนรู้การบีบอัดและคลายข้อมูล โดยข้อมูลที่มีค่า Reconstruction Error สูงผิดปกติ

\begin{myBox}[]{วิธีการตรวจจับความผิดปกติด้วย Z-Score และ IQR}{}
สมมติมีข้อมูลเดียวมิติ (Univariate Data) $x_1, x_2, \ldots, x_N$ วิธีการตรวจจับความผิดปกติมีดังนี้

\textbf{Z-Score Method:} คำนวณ Z-Score ของแต่ละข้อมูล
$$z_i = \frac{x_i - \bar{x}}{s}$$
โดยที่ $\bar{x}$ คือค่าเฉลี่ยและ $s$ คือส่วนเบี่ยงเบนมาตรฐาน ข้อมูลที่มี $|z_i| > 3$ (ห่างจากค่าเฉลี่ยเกิน 3 ค่าเบี่ยงเบนมาตรฐาน) ถือเป็นความผิดปกติ

\textbf{IQR Method:} คำนวณ Quartiles และ IQR
\begin{align*}
Q_1 &= 25\text{th percentile} \\
Q_3 &= 75\text{th percentile} \\
\text{IQR} &= Q_3 - Q_1
\end{align*}
ขอบเขตความผิดปกติ (Anomaly Boundaries):
$$[\,Q_1 - 1.5 \times \text{IQR},\; Q_3 + 1.5 \times \text{IQR}\,]$$
ข้อมูลที่อยู่นอกช่วงนี้ถือเป็นความผิดปกติ
\end{myBox}

ในบริบทของข้อมูลหลายมิติ (Multivariate Data) การตรวจจับความผิดปกติมีความซับซ้อนมากขึ้น เนื่องจากต้องพิจารณาความสัมพันธ์ระหว่างหลายตัวแปรพร้อมกัน วิธี Mahalanobis Distance เป็นวิธีที่นิยมใช้สำหรับข้อมูลหลายมิติ โดยวัดระยะทางจากจุดข้อมูลถึงศูนย์กลางของการแจกแจง โดยคำนัญถึงโครงสร้างความแปรปรวนร่วม (Covariance Structure) ของข้อมูล

\begin{myBox}[]{Mahalanobis Distance สำหรับข้อมูลหลายมิติ}{}
สำหรับข้อมูลหลายมิติ $\mathbf{x} \in \mathbb{R}^d$ ที่มี Covariance Matrix $\mathbf{S}$ Mahalanobis Distance จาก $\mathbf{x}$ ถึงค่าเฉลี่ย $\boldsymbol{\mu}$ คำนวณโดย

$$D_M(\mathbf{x}) = \sqrt{(\mathbf{x} - \boldsymbol{\mu})^\top \mathbf{S}^{-1} (\mathbf{x} - \boldsymbol{\mu})}$$

Mahalanobis Distance จะเท่ากับ Euclidean Distance เมื่อ $\mathbf{S}$ เป็นเมทริกซ์เอกลัพธ์ (Identity Matrix) ซึ่งหมายความว่าตัวแปรทั้งหมดมีความแปรปรวนเท่ากันและไม่มีความสัมพันธ์กัน

ข้อมูลที่มี $D_M(\mathbf{x}) > \chi^2_{d, 0.975}$ (ค่า Critical ของ Chi-Square Distribution ที่ $d$ มิติ ที่ระดับนัยสำคัญ 0.05) ถือเป็นความผิดปกติ
\end{myBox}

การตรวจจับความผิดปกติมีความท้าทายเฉพาะตัวที่แตกต่างจากการจำแนกประเภทหรือการถดถอย ประการแรก ความถ่วงแบบเบ้ (Class Imbalance) รุนแรงมาก เนื่องจากความผิดปกติมักมีจำนวนน้อยมากเมื่อเทียบกับข้อมูลปกติ (เช่น ธุรกรรมฉ้อโกง 1 ราย ในธุรกรรมปกติ 10,000 ราย) ประการที่สอง ไม่มีค่า Ground Truth ที่แน่นอน เนื่องจากความผิดปกติอาจเปลี่ยนแปลงรูปแบบตามเวลา (Concept Drift) ประการที่สาม ข้อมูลที่ถูกติดป้ายว่าเป็นความผิดปกติอาจมีความหลากหลายมาก ทำให้ยากที่จะสร้างตัวแบบที่ครอบคลุมทุกรูปแบบ

\begin{exmp}
เพื่อทำความเข้าใจการตรวจจับความผิดปกติอย่างเป็นรูปธรรม พิจารณาตัวอย่างการตรวจจับธุรกรรมทางการเงินที่ผิดปกติ สมมติมีข้อมูลธุรกรรม 10,000 รายการ โดยแต่ละธุรกรรมมียอดเงิน (Amount) และเวลาที่ทำธุรกรรม (Hour of Day) เป็นคุณลักษณะ ธุรกรรมส่วนใหญ่ (ปกติ) มียอดเงินไม่เกิน 10,000 บาท และทำธุรกรรมในช่วงเวลาทำการ (8.00-18.00 น.) แต่มีธุรกรรมผิดปกติ เช่น ธุรกรรมวงเงินสูงมากในเวลาวิกาล หรือธุรกรรมที่มีรูปแบบแตกต่างจากพฤติกรรมปกติของลูกค้า

ในการตรวจจับด้วย Mahalanobis Distance จะคำนวณค่าเฉลี่ย $\boldsymbol{\mu} = (\bar{x}_{amount}, \bar{x}_{hour})$ และ Covariance Matrix $\mathbf{S}$ จากธุรกรรมที่เป็นปกติ จากนั้นสำหรับธุรกรรมใหม่แต่ละรายการ จะคำนวณ $D_M$ และถ้าค่าเกินค่า Threshold ที่กำหนด (เช่น $\chi^2_{2, 0.975} \approx 7.38$) ระบบจะแจ้งเตือนว่าเป็นธุรกรรมที่น่าสงสัย
\end{exmp}

การประยุกต์ใช้งานการตรวจจับความผิดปกติมีความกว้างขวางมาก ในด้านความปลอดภัยทางไซเบอร์ (Cybersecurity) ใช้ในการตรวจจับการบุกรุกระบบ (Intrusion Detection) การโจมตีแบบ DDoS และมัลแวร์ ในด้านการแพทย์ (Medical Diagnosis) ใช้ในการตรวจจับเซลล์มะเร็งจากผลตรวจที่ผิดปกติ การตรวจจับภาวะหัวใจเต้นผิดจังหวะ (Arrhythmia) จากสัญญาณ ECG และการระบุผู้ป่วยที่มีความเสี่ยงสูงจากผลตรวจทางห้องปฏิบัติการ ในด้านอุตสาหกรรม (Industrial Fault Detection) ใช้ในการตรวจจับความผิดปกติของเครื่องจักรจากข้อมูลเซ็นเซอร์ การตรวจจับการรั่วไหลของสารเคมีจากการเปลี่ยนแปลงของเซ็นเซอร์ และการเฝ้าระวังคุณภาพผลิตภัณฑ์ในสายการผลิต

\begin{Verbatim}[tabsize=4]
# Anomaly Detection in Python
import numpy as np
from scipy import stats
import matplotlib.pyplot as plt

# Sample data: 1000 normal transactions (amount, hour)
np.random.seed(42)
normal_data = np.random.normal(loc=[5000, 10], scale=[2000, 3], size=(1000, 2))

# Add 5 anomalies
anomalies = np.array([[50000, 3], [45000, 2], [48000, 4], [52000, 3], [49000, 2]])
data = np.vstack([normal_data, anomalies])

# Z-Score method for each feature
z_scores = np.abs(stats.zscore(data, axis=0))
anomalies_zscore = np.where((z_scores > 3).any(axis=1))[0]

# Mahalanobis Distance method
mean = np.mean(data[:-5], axis=0)  # mean of normal data only
cov = np.cov(data[:-5], rowvar=False)
cov_inv = np.linalg.inv(cov)

def mahalanobis(x, mean, cov_inv):
    diff = x - mean
    return np.sqrt(np.dot(np.dot(diff, cov_inv), diff))

mahal_dist = np.array([mahalanobis(x, mean, cov_inv) for x in data])
threshold = stats.chi2.ppf(0.975, df=2)  # 7.38 for 2D
anomalies_mahal = np.where(mahal_dist > threshold)[0]

print(f"Z-Score detected {len(anomalies_zscore)} anomalies")
print(f"Mahalanobis detected {len(anomalies_mahal)} anomalies")

# Plot results
plt.scatter(data[:-5, 0], data[:-5, 1], label='Normal', alpha=0.5)
plt.scatter(data[anomalies_mahal, 0], data[anomalies_mahal, 1], 
            color='red', marker='x', s=100, label='Anomaly')
plt.xlabel('Transaction Amount (THB)')
plt.ylabel('Hour of Day')
plt.legend()
plt.title('Anomaly Detection using Mahalanobis Distance')
plt.show()
\end{Verbatim}

ข้อจำกัดและความท้าทายของการตรวจจับความผิดปกติที่ควรตระหนักมีหลายประการ ประการแรก การตั้ง Threshold ที่เหมาะสมเป็นเรื่องยาก เนื่องจากถ้าตั้งสูงเกินจะพลาดความผิดปกติจริง ๆ (False Negative สูง) แต่ถ้าตั้งต่ำเกินจะแจ้งเตือนผิดพลาดมากเกินไป (False Positive สูง) ประการที่สอง ความผิดปกติที่ซับซ้อนอาจไม่สามารถตรวจจับได้ด้วยวิธีทางสถิติเ� alone ต้องใช้เทคนิคที่ซับซ้อนกว่าเช่น Deep Learning ประการที่สาม ความผิดปกติอาจเปลี่ยนแปลงรูปแบบตามเวลา (Concept Drift) ทำให้ตัวแบบที่เคยใช้ได้อาจล้าสมัยอย่างรวดเร็ว

\newpage
%\RefChapter
